\scanpage{24}
\mspage{22}
l'isomorphisme garantira automatiquement qu'il est compatible avec $F$ et $V$).

2.5. Je n'ai pas vérifié, à vrai dire, les deux affirmations, mais n'ai pas le moindre doute qu'elles sont correctes telles quelles. Cette façon de voir le module de Dieudonné permet de plus d\deleted{e}'expliciter de façon remarquable les variations infinitésimales de structure de la variété abélienne $A$, donnée\inserted{,} en caractéristique \deleted{nulle} $p>0$, en termes du module de Dieudonné: les prolongements de $A$ en un schéma abélien $B$ sur $\Lambda$ doivent correspondre exactement aux \deleted{sous} modules quotients libres de rang $g$ de $M\otimes_W\Lambda$ qui redonnent, modulo $p$, le module quotient de rang $g$ canonique de $M\otimes_Wk=H^1_{\operatorname{DR}}(A)$, (correspondant à la filtration canonique de cette cohomologie). Plus généralement et plus précisément:

3º) Considérons, pour tout schéma abélien $A$ sur une base $S$, sur la cohomologie de De Rham $R^1f_*(\underline{\Omega}^{\bullet}_{A/S})=H^1(f)$, la filtration canonique
\[
0\longleftarrow R^1f_*(\underline{O}_A)\longleftarrow H^1(f)\longleftarrow R^0f_*(\underline{\Omega}^1_{A/S})\longleftarrow0.
\]
Donc pour tout $B$ sur $T$ comme dans 1º, on a sur \inserted{$\underline{M}(T)=$} $\operatorname{DR}^1(A/S)(T)=H^1(g)$ une filtration naturelle, ne dépendant que de la classe à isomorphisme près du relèvement choisi $B$ de $A$ à $T$, et prolongeant la filtration déjà connue de \deleted{$\operatorname{DR}^1($} $\underline{M}(S)=H^1(f)$ provenant de $A$. Ceci dit, \uline{troisième affirmation}: \uline{on obtient ainsi une correspondance biunivoque entre classes de prolongements de} $A$ \uline{à} $T$, \uline{et prolongements de la filtration donnée de} $\underline{M}(S)=\underline{M}(T)\otimes_{\underline{O}_T}\underline{O}_S$ \uline{en une filtration de} $\underline{M}(T)$. Plus précisément encore, le foncteur $B\rightsquigarrow(A,\phi)$, de la catégorie des schémas abéliens \inserted{$B$} sur $T$, dans la catégorie des schémas abéliens \deleted{sur} $A$ sur $S$, munis d'une filtration \inserted{$\phi$} de $\operatorname{DR}^1(A/S)(T)$ prolongeant celle de $\operatorname{DR}^1(A/A)(S)$ (NB il ne s'agit que de filtrations à quotients localement libres, bien sûr) est une équivalence de catégories.

Ce\deleted{s}\inserted{t} énoncé\deleted{s} est évidemment fort suggestif aussi pour des généralisations en cohomologie de De Rham de dimension supérieure, pour un morphisme lisse et projectif quelconque, tenant compte de la filtration canonique de celle ci. On voit bien en tous cas que ce dernier élément de structure n'est \uline{pas} de nature cristalline, i.e. donnée par une filtration du cristal de De Rham postulé dans 2.2., mais est au contraire dans la nature d'un élément de structure ``conti-
\margin{23}
